\chapter{¿Es esto un dispositivo?}
\label{cap:prospectiva}

Este libro se llama \textit{El dispositivo caldereño} y nunca definió la palabra. La usó decenas de veces, la puso en la tapa y la heredó del volumen anterior, pero no la sometió a la
misma exigencia que le impuso a cada expediente.

\textbf{Este capítulo hace eso, y con el mismo criterio que el resto: el término se conserva si
lo documentado lo sostiene, y se abandona si no.} De esa respuesta depende, además, la única
pregunta que el libro no ha contestado: qué va a pasar acá.

\begin{cautela}
\textbf{Es el único capítulo del libro que proyecta, y el único que argumenta más de lo que documenta} ---el capítulo \ref{cap:infraestructura} también relee, pero para describir, y el \ref{cap:plan} propone, pero no pronostica---. Casi no agrega hechos nuevos: reordena los que los veintiséis capítulos anteriores establecieron, y la única novedad ---el lanzamiento del PDES 2050--- se consigna como tal. Quien lo lea salteado debe saber que \textbf{acá se razona, y que en
el resto se transcribe}.
\end{cautela}

\section{Qué quiere decir la palabra}

El término viene de Michel Foucault, que lo usó en 1977 para nombrar una clase de objeto que la
palabra «institución» no alcanza a describir. Un dispositivo es, en su formulación, un
\textbf{conjunto deliberadamente heterogéneo}: discursos, instituciones, disposiciones
arquitectónicas, decisiones reglamentarias, leyes, medidas administrativas, enunciados
científicos. Lo que lo constituye no es ninguno de esos elementos sino \textbf{la red que puede
establecerse entre ellos}. Y tiene una propiedad decisiva: cumple una \textbf{función
estratégica} y responde a una urgencia, \textbf{sin que haya un estratega}.

Giorgio Agamben amplió el término en 2006 hasta volverlo casi ilimitado: dispositivo es
cualquier cosa que tenga la capacidad de capturar, orientar, determinar, interceptar, modelar,
controlar o asegurar las conductas de los seres vivos.

\begin{cautela}
\textbf{La versión de Agamben no sirve para este libro, y conviene decir por qué.} Si
dispositivo es cualquier cosa que orienta conductas, entonces todo lo es ---una ruta, una tarifa,
un horario--- y el término deja de distinguir nada. \textbf{Un concepto que no excluye nada no
prueba nada.}

Este capítulo usa la versión estricta: \textbf{red heterogénea, función estratégica, ausencia de
estratega}. Es más difícil de sostener y por eso vale más.
\end{cautela}

\section{Primera condición: la heterogeneidad}

\textbf{Se cumple, y con holgura.} Lo que este libro documenta sobre un territorio de dos
municipios incluye, en la misma trama:

Un poseedor colonial del siglo XVII, partidos policiales de 1909 y sucesiones del siglo XX \textit{(cap.\ \ref{cap:fincas})}. Una ley
de 1929 que autoriza a un municipio de otro departamento a reconcentrar el agua de esta cuenca
\textit{(cap.\ \ref{cap:aguabaja})}. Una expropiación de 1972 para un embalse que abastece a la
capital \textit{(cap.\ \ref{cap:expropiacion})}. Un código de aguas de 1998, reglamentado en 2003, con un decreto
de línea de ribera de 2002 \textit{(cap.\ \ref{cap:ribera})}. Un plan de desarrollo aprobado con carácter definitorio \textit{(cap.\ \ref{cap:pdua})}. Un código de ordenamiento municipal que declara no urbanizables las márgenes del río sin decir cuántos metros \textit{(cap.\ \ref{cap:cot})}. Veinte
concesiones de áridos ante un juzgado de minas \textit{(cap.\ \ref{cap:resistencias})}. Un
ordenamiento de bosques nativos con su ley, su cartografía y su umbral de deforestación por
cuenca \textit{(cap.\ \ref{cap:opacidad})}. Un amparo colectivo \textit{(cap.\
\ref{cap:amparo})}. Una obra vial con audiencia pública \textit{(cap.\ \ref{cap:vaqueros})}.
Contratos de promoción turística con exención de tributos \textit{(cap.\ \ref{cap:poblacion})}. Y
una gacetilla vecinal de 2018 \textit{(cap.\ \ref{cap:resistencias})}.

\textbf{Nada de eso pertenece al mismo campo, ni a la misma época, ni a la misma autoridad.} Y
sin embargo recae todo sobre el mismo suelo, y este libro pudo reconstruir la red que los une
sin forzar un solo dato.

\section{Segunda condición: la función estratégica}

\textbf{También se cumple, y este libro puede enunciarla en una frase.}

\begin{ficha}{La función, dicha con los números del libro}
El departamento \textbf{produce agua} para una ciudad que no es la suya: el embalse de Campo
Alegre con su planta de cuatro mil metros cúbicos por hora, el Acueducto Norte, la ley de 1929.

\textbf{Produce material pétreo} para esa misma ciudad: veinte concesiones sobre los cauces, y
dos de las concesionarias son contratistas de obra pública de la capital \textit{(cap.\
\ref{cap:infraestructura})}.

\textbf{Produce suelo residencial} para quienes trabajan allá: crecimiento poblacional del
58,4 por ciento en doce años, la tasa de empleo más alta de la provincia, y ---en el municipio de La Caldera--- \textbf{2,23
personas por vivienda}, el tercer valor más bajo entre sesenta municipios; en el departamento, 2,45 contra 3,13 provinciales.

\textbf{Y no recibe lo equivalente}: entre el dieciocho y el cincuenta y ocho por ciento de sus
hogares sin agua de red según la zona, más del noventa y cinco por ciento sin cloaca, y una
calzada de cinco metros y medio que no puede ensancharse.
\end{ficha}

\textbf{Tres flujos en la misma dirección, sostenidos por instrumentos de tres siglos y de
cuatro jurisdicciones distintas.} Eso es una función estratégica en el sentido preciso del
término: no una intención, sino un efecto sistemático que los elementos producen por su
disposición relativa.

\section{Tercera condición: que no haya estratega}

\textbf{Esta es la que el libro tardó más en poder afirmar, y es la que decide todo.}

\begin{ficha}{Lo que se buscó y no se encontró}
\textbf{No hay un plan que lo decida.} El PDES 2030 se aprobó con carácter definitorio invocando el artículo 77 de la Constitución provincial, y «La Caldera» aparece \textbf{dos veces en ciento setenta páginas}: en la nómina de
asistentes al Diagnóstico y en un inventario hotelero. Ese mismo plan cartografía los salares de
la puna con sus empresas y sus toneladas: \textbf{sabe territorializar cuando quiere}.

\textbf{No hay una decisión de perfil vigente.} Hubo actos que intentaron asignarle una ---la Ley 1423 (orig.\ 145) de 1934 incorporó al departamento a la región económica del olivo (capítulo \ref{cap:hacienda}) y el Decreto 12 de 1980 reconoció su área tabacalera (capítulo \ref{cap:tierra})---, pero ninguno de los dos perfiles sobrevivió, y no hay acto posterior que asigne al corredor vocación turística, agropecuaria ni residencial.

\textbf{Y el Estado lo reconoce por escrito.} El propio plan declara, en 2012, que el
crecimiento del área metropolitana fluyó «\textit{espontáneamente, \textbf{sin
planificación}}», y que hay que «\textit{paliar los efectos de la falta de previsión}».
\end{ficha}

\textbf{No hay estratega.} Y no porque este libro no lo haya encontrado, sino porque el
instrumento que debería contenerlo declara que no existe.

\textbf{Y hay una prueba más, de otra clase, que los barridos de 1918 y 1919 pusieron a la vista.}
En marzo de 1918 un gobierno decidió, por un umbral de \textbf{renta}, que catorce departamentos
pasaban a ser Municipalidades electivas, y a La Caldera se lo revocó rehaciendo la cuenta. En
enero de 1919 el gobierno siguiente declaró que ese umbral no existía: que el artículo en que se
fundaba estaba derogado y que lo que la Constitución exige es un umbral de \textbf{población}
(capítulo \ref{cap:siglo}). \textbf{Doce meses, dos gobiernos, dos reglas distintas para la misma
pregunta, y ninguno de los dos citando al otro.} Un dispositivo con estratega tiene una regla; lo
que el archivo muestra en estos dos años es que \textbf{ni siquiera la regla estaba establecida},
y que el resultado ---el departamento sin gobierno propio--- salió igual por los dos caminos.

\textbf{Y el barrido de 1922 muestra hasta dónde llegó ese resultado.} El 18 de abril de ese año
un Interventor Nacional designa las Comisiones Municipales de los departamentos «\textit{que en
la actualidad carecen de ella}», y La Caldera va primera en la lista (capítulo
\ref{cap:siglo}). \textbf{No es que el departamento tuviera un gobierno débil: no tenía ninguno},
y el archivo no registra cuándo lo perdió. Cuatro años después de la discusión sobre si debía ser
Municipalidad electiva, el cuerpo que se le discutía había desaparecido sin que ningún acto
publicado lo dijera.

\textbf{Y diecisiete meses después la Provincia intervino el que acababa de crearse.} El decreto
1161 del 24 de septiembre de 1923 declara intervenida la Municipalidad del departamento y nombra
un Interventor, a pedido de presentaciones que un expediente instruye y que no se publican
(capítulo \ref{cap:siglo}). \textbf{Con eso, el ciclo se cierra y se puede contar entero}: en
1918 el Estado resolvió que el departamento no estaba en condiciones de tener municipalidad
electiva, por un umbral de renta; en 1919 declaró que el umbral era la población; entre julio de
1921 y abril de 1922 el departamento se quedó sin Comisión Municipal sin que ningún acto lo
dijera; en abril de 1922 un Interventor Nacional se la devolvió; en septiembre de 1923 el
gobierno provincial la intervino; \textbf{y en abril de 1924, aceptadas las renuncias
«indeclinables» de dos de los miembros que quedaban, el Poder Ejecutivo designó por decreto a los
cuatro que la integrarían} ---uno de ellos el mismo que había renunciado tres meses antes, y otro
el concesionario del agua--- \textbf{«vistas las comunicaciones» de un expediente que no se
publica}. \textbf{Seis actos en siete años, ninguno de los cuales cita al
anterior, y los seis resuelven quién administra este territorio. Ninguno lo decide el
territorio}, y en ninguno de los seis hay una elección.

\textbf{Y la serie no termina ahí: llega hasta 1946 y suma dos episodios más.} En septiembre de
1930 los actos del departamento pasan a firmarlos un \textbf{Interventor Nacional}; y entre 1944
y agosto de 1946 el cuerpo municipal está otra vez \textbf{intervenido}, esta vez sin que ningún
acto publicado diga cuándo se dispuso (capítulos \ref{cap:siglo} y \ref{cap:opacidad}).
\textbf{Ocho actos en veintiocho años, y de los ocho ninguno sale de una elección.} La lámina
\ref{fig:cargos} los marca con línea vertical: 1918, 1922, 1923, 1930 y 1944.

\textbf{Y hay un noveno hecho, que no es un acto sino su ausencia}: entre julio de 1921 y abril
de 1922 el departamento estuvo \textbf{sin comisión municipal}, y el archivo no registra cuándo
la perdió. \textbf{De veintiocho años, este libro puede documentar al menos nueve en que el
cuerpo municipal no era un cuerpo electo}: los de las dos designaciones por interventor nacional,
los de las dos intervenciones y el año sin comisión.

\textbf{Ésa es la forma más pura de la tercera condición}: no hay estratega
porque, en los tramos donde el archivo deja mirar, no hay siquiera un órgano local estable al que
un estratega pudiera dirigirse.

\textbf{Y el arco del agua da la prueba positiva que a esta condición le faltaba.} Entre 1885 y
1946 el municipio pasa de \textbf{conceder} treinta litros por segundo del Wierna a
\textbf{prestar «la colaboración necesaria»} a la Dirección General de Hidráulica, en seis actos
de los cuales \textbf{ninguno deroga al anterior} (lámina \ref{fig:agua} y capítulo
\ref{cap:infraestructura}). \textbf{Eso es exactamente lo que este capítulo entiende por
dispositivo y no por plan}: nadie decidió quitarle la facultad, y sin embargo al final no la
tiene. \textbf{Si hubiera habido un estratega, habría un acto que la deroga}; lo que hay es una
sucesión de actos que la rodean.

\textbf{Y hay una razón anterior, que el archivo temprano dejó ver y que vale más que la
anterior porque no depende de lo que un plan declare de sí mismo: en este departamento nunca
hizo falta un estratega, porque el conjunto de personas que podían ocupar los cargos era
demasiado chico.} En 1909, con ciento setenta y siete inscriptos en el registro cívico, el
hombre que clasificaba a los contribuyentes para el impuesto de patentes pasó a mandar la
policía del departamento; el juez de paz suplente era, al mismo tiempo, comisario suplente y
miembro de la comisión municipal; el comisario auxiliar de un partido era propietario en ese
partido; y el presidente municipal que firmaba el cuadro de rentas administraba la finca más
grande del territorio, que su propia familia tenía en condominio (capítulos \ref{cap:fincas} y
\ref{cap:siglo}). Los tres reemplazos que la comisión municipal necesitó ese año salieron del
mismo jurado de patentes del año anterior.

Eso no prueba acuerdo ni parentesco, y este libro no lo afirma. Prueba algo distinto y más
duradero: que \textbf{la coincidencia entre quien regula y quien es regulado no necesitó
producirse, porque venía dada}. Seis años más tarde vuelve a ocurrir sin que nadie lo decida: en diciembre de 1915 el juez de paz del departamento es el comisionado que clasifica sus patentes, y los dos jurados que resolverán los reclamos contra esa clasificación son dos de sus propietarios mayores, con un miembro de la comisión municipal de suplente (capítulo \ref{cap:siglo}). Un siglo después, el capítulo \ref{cap:defensas} encuentra la misma
forma con otros nombres: la sociedad que en un mismo año es concesionaria de áridos sobre tierra
fiscal, oferente admisible en la obra que protege el acueducto y contratista directa de su
reparación. \textbf{Y una tercera vez en 1918}, ya no en el pueblo chico de 1909 sino bajo una Intervención
Nacional: el presidente de la Mesa N.\ 1 del Colegio Electoral del departamento es, en el mismo
año, el juez de agua de uno de sus dos distritos, nombrado por el Interventor a propuesta del
propio municipio; su segundo suplente de mesa es el mismo Augusto Regis\index{Regis, Augusto}
que el archivo registra desde 1908, y el apellido que los deslindes de la Calderilla de 1917
registran como propietario lindero figura en la mesa de Mojotoro como autoridad electoral
(capítulo \ref{cap:siglo}). \textbf{Tres veces en nueve años, con tres juegos de nombres y bajo
tres regímenes distintos ---gobierno electo, gobierno electo otra vez, intervención federal---,
y ninguna de las tres necesitó que nadie la arreglara. \textbf{Y una cuarta en 1919}, con el
Presidente del Senado en ejercicio del Ejecutivo.} La condición se sostiene, entonces, por dos vías
independientes: el instrumento de 2012 que declara que nadie planificó, y el archivo de 1909 que
muestra por qué no hacía falta que nadie lo hiciera ---una serie, no un episodio: 1909,
1915, 1918, 1919, 1920 y 1921, con seis juegos de nombres y bajo cinco gobiernos distintos. En 1919, dos
de los cinco miembros de la Comisión Municipal son, en el mismo trimestre, comisarios de partido
bajo las órdenes del Jefe de Policía, y ese mismo cuerpo eleva las ternas de jueces de paz; y el
veedor del Código Rural de Chalcanio de 1909 es, diez años después, el comisario auxiliar de ese
mismo paraje (capítulo \ref{cap:siglo}). \textbf{Y en 1920 la superposición cambia de especie},
que es lo que la vuelve interesante: ya no es un cargo que coincide con otro cargo, sino que
\textbf{uno de los tres candidatos a diputado por el departamento es quien dos años antes había
pedido las dos únicas concesiones mineras que el Boletín registra sobre este suelo}. Coincidencia
de nombre y no identificación registral, y aun así es el primer caso de la serie en que quien
pide una concesión sobre el departamento se presenta a representarlo. \textbf{Y en 1921 vuelve a
cambiar de especie}: \textbf{Daniel Linares}\index{Linares, Daniel} es, el mismo año, primer
suplente de la mesa 1 del departamento y el propietario que dos de los cuatro rumbos del deslinde
de la finca El Angosto declaran como lindero (capítulo \ref{cap:siglo}). \textbf{Y en 1922 una séptima, bajo otra Intervención Nacional y con la forma más neta de
todas}: \textbf{Jorge Rauch} es comisario de policía del departamento desde marzo y miembro de su
Comisión Municipal desde abril; \textbf{Carlos F.\ Matienzo}\index{Matienzo, Carlos} acumula, con quince días de
diferencia, el juzgado de paz propietario y el Registro Civil del mismo departamento; y
\textbf{Daniel Linares}\index{Linares, Daniel}, que en enero recibe la correspondencia electoral del departamento, en
abril integra el cuerpo que lo gobierna (capítulo \ref{cap:siglo}). \textbf{Y en 1923 una octava, que cambia de escala}: el escrito con que
Francisco S.\ Urquiza\index{Urquiza, Francisco S.} pide la primera concesión de agua del
departamento declara que su acequia atraviesa la finca Vaqueros, «\textit{de propiedad de los
herederos Toranzos y doctor \textbf{Carlos Serrey}}»\index{Serrey, Carlos}, y cinco meses después
ese mismo Carlos Serrey se incorpora al Senado provincial por Cerrillos, en la misma nómina que
trae al senador por Caldera (capítulo \ref{cap:siglo}). \textbf{Ya no es un cargo que coincide con
otro cargo dentro del departamento: es el propietario del suelo que la obra debe atravesar
sentado en la legislatura de la provincia que concede el agua.} Coincidencia de nombre y no
identificación registral. \textbf{Y en 1924 una novena, que es la más completa de todas.} El mismo Urquiza\index{Urquiza, Francisco S.}
aparece cuatro veces en un año y en cuatro calidades: \textbf{en enero} recibe la concesión de
agua, fundada en el informe de la Comisión Municipal del departamento; \textbf{en marzo} figura
como propietario lindero en el deslinde de la finca vecina; \textbf{en abril} el decreto 1576 lo
designa \textbf{miembro de esa misma Comisión Municipal}; y \textbf{en octubre} preside la mesa de
Vaqueros (capítulo \ref{cap:siglo}). \textbf{Tres meses separan al beneficiario de la primera
concesión de agua del departamento del cuerpo cuyo informe la fundó}, y veinte años después el
mismo hombre es el mayor de los cuarenta y siete titulares que publica el revalúo de 1944.
\textbf{Y el barrido de 1925 agrega una quinta calidad}: el decreto 2674 del 25 de junio de ese
año nombra subcomisario \textit{ad honorem} de Vaqueros a Antonio Agüero «\textit{en reemplazo de
D.\ Francisco S.\ Urquiza}», de modo que el mismo hombre era también \textbf{la autoridad policial
del paraje}, sin sueldo y sin que su nombramiento aparezca en ningún tramo leído.
\textbf{Y en 1940 una décima, que es la más directa de todas}: el Boletín
identifica al solicitante de una ayuda para las fiestas patronales del pueblo
como «\textit{el señor presidente de la comisión municipal del distrito de la
Caldera, \textbf{diputado provincial don Augusto Regis}}»\index{Regis, Augusto}
(capítulo \ref{cap:siglo}). \textbf{No es un cargo que coincide con otro dentro
del departamento: es la misma persona presidiendo el municipio y sentada en la
Legislatura que lo regula}, y lleva para entonces cuatro reelecciones seguidas
por decreto y dieciocho años entrando y saliendo de ese cuerpo. \textbf{Y en 1942 una undécima, que completa el cuadro}: el mismo Augusto
Regis\index{Regis, Augusto} es nombrado \textbf{comisario de policía de La
Caldera} mientras preside la Comisión Municipal del distrito ---reelecto por
decreto ese año y el siguiente--- y ocupa una banca de diputado provincial
(capítulo \ref{cap:siglo}). \textbf{Presidente del cuerpo que gobierna el
departamento, comisario de su policía y legislador de la provincia que lo
regula, a la vez y en actos publicados.} \textbf{Y la relectura de 1911 y 1912 agrega una que es
anterior a todas menos la primera, y de otra especie}: en octubre de 1911 la Provincia crea, \textit{ad
honorem}, una oficina del Registro Civil en la estación Mojotoro y la pone a cargo de Benjamín
López\index{López, Benjamín}, subcomisario del partido de la estación; en diciembre de 1912 el
mismo López pasa a comisario del partido, y en el juicio de ese año declara como jefe de la
estación del ferrocarril (capítulo \ref{cap:presencia}). \textbf{El registro de las personas, la
policía y el ferrocarril, en un solo hombre, en el mismo paraje y sin sueldo por el primero.}
\textbf{Doce ocurrencias en treinta y tres años, y ya no todas del mismo tipo}: cargo con cargo, concesionario con
candidato, autoridad electoral con propietario del suelo que se deslinda, dos ramas del
Estado local ---la justicia de paz y el registro de las personas--- en la misma persona y en la
misma quincena, propietario afectado con legislador provincial, concesionario de agua
con miembro del cuerpo que informó su concesión, y tres funciones del Estado en un solo hombre
allí donde llegaba el tren.

\begin{cautela}
\textbf{Conviene distinguir esto de dos cosas que no son.}

\textbf{No es una conspiración.} Nadie se reunió a decidir que La Caldera proveyera agua, piedra
y dormitorios. Los actos que producen ese resultado tienen cada uno su fundamento propio y su
autoridad competente, y este libro no encontró uno solo que sea ilegal.

\textbf{Y no es un accidente.} Un accidente no se repite durante tres siglos en la misma
dirección. Lo que hace que los tres flujos vayan siempre hacia el mismo lado no es el azar: es
que \textbf{los instrumentos están dispuestos de manera tal que ése es el resultado que
producen}, y nadie tiene a su cargo mirar el conjunto.
\end{cautela}

\section{Entonces sí, y el término dice algo}

Las tres condiciones se cumplen. \textbf{La Caldera es un dispositivo en el sentido estricto: una
red heterogénea de instrumentos que cumple una función estratégica sin que nadie la haya
diseñado.}

\textbf{Y decirlo así no es una figura literaria: tiene consecuencias prácticas.}

Si hubiera un diseño, la corrección sería política: cambiar la decisión. \textbf{Como no lo hay,
la corrección es institucional}: producir el órgano, el registro o el criterio que hoy no
existen. Por eso el apéndice \ref{cap:propuestas} contiene veintiséis propuestas normativas y
ninguna denuncia.

Si hubiera un culpable, el libro sería una acusación. \textbf{Como no lo hay, es un inventario}:
doscientos setenta y ocho pedidos, ordenados en doce secciones ---once por destinatario y una sobre la vida cotidiana---, y la mayoría pide lo que la propia norma
manda producir.

\section{Qué viene, entonces}

\begin{cautela}
\textbf{Todo lo que sigue es prospectivo.} No documenta: proyecta sobre lo documentado.
Cualquiera de estas líneas puede no ocurrir, y ninguna debe citarse como hallazgo de este libro.
\end{cautela}

\subsection*{Lo que la dinámica indica}

Si no hay diseño, lo que viene no lo decide nadie: lo produce la disposición de los elementos.
Y esa disposición apunta en una dirección \textbf{que no es la que el propio departamento se
atribuye}.

\textbf{No es turismo.} Hoteles y restaurantes ocupan a doscientas dieciocho personas, el 3,49
por ciento: por debajo del promedio provincial y por debajo de la capital. El servicio doméstico
emplea al doble. Las veinte canteras, a cincuenta y una personas.

\textbf{No es agricultura.} Ciento noventa y tres personas, sobre una superficie declarada que
cayó un ochenta y ocho por ciento ---aunque el capítulo \ref{cap:tierra} muestra que lo que cayó
fue sobre todo declaración de monte, no de cultivo---.

\textbf{Es residencia.} Población en alza sostenida, empleo alto, viviendas semivacías, el Estado como principal empleador y una parte del trabajo, presumiblemente, en la capital. \textbf{Lo que los números describen es un suburbio residencial y de fin de semana de una
capital de medio millón de habitantes.}

\subsection*{Lo que lo acelera}

\textbf{La ruta.} La pavimentación del tramo entre la rotonda de avenida Bolivia y el puente del
río La Caldera, expediente 0110033-243466/2021-181, con estudio de impacto y audiencia convocada para noviembre de 2023. El capítulo \ref{cap:vaqueros} establece el mecanismo: la infraestructura que reduce
la fricción con la capital valoriza el suelo.

\textbf{Y no es nuevo.} Entre 1923 y 1945 lo que el Estado construyó en el departamento fueron
caminos, y los caminos seguían de largo, hacia Campo Santo, Güemes, Jujuy y la capital (capítulo
\ref{cap:presencia}). La ruta que hoy acelera la valorización es la capa más reciente de la
inversión que el Estado sostuvo sin pausa en ese tramo del archivo.

\textbf{El agua.} La red se está extendiendo ---barrio Santiago Apóstol, septiembre de 2026--- y
el capítulo \ref{cap:redes} muestra que la factibilidad hídrica es la condición práctica del
fraccionamiento.

\subsection*{Lo que lo limita, y es lo que este libro no esperaba}

\textbf{La ladera está ocupada.} El titular provincial de Vialidad Nacional declara en agosto de
2026 que la calzada tiene cinco metros y medio, y la nota de prensa que recoge esa declaración
agrega, sin atribuírselo entre comillas, que \textbf{los emprendimientos privados sobre
las laderas, muchos sin autorización, complican cualquier ampliación futura del corredor}
(capítulo \ref{cap:vaqueros}).

\textbf{La pendiente está protegida.} El nuevo artículo 14 de la Ley 7543 pone en Categoría II
todo lo que tiene \textbf{pendiente superior al quince por ciento}, que es exactamente esa
ladera.

\textbf{Y desde diciembre de 2024 hay un permiso más.} El artículo 20 reformado exige que toda
urbanización con cambio de uso de suelo en zona categorizada cuente con \textbf{la no objeción
de la autoridad forestal}.

\textbf{Y el bosque de esta cuenca no es categorizable.} Es el límite más duro de los cuatro y
el que este capítulo no había puesto en esta lista, aunque el libro lo documenta: el capítulo
\ref{cap:opacidad} establece, con las capas del propio informe técnico de la Ley 8483, que la
cuenca Mojotoro--Lavayén ---que ocupa el setenta y nueve por ciento del departamento--- figura
excedida en su umbral máximo de deforestación y tiene \textbf{cero hectáreas categorizables a
verde}. Según el criterio que el mismo informe enuncia, en una cuenca excedida no se autorizan
proyectos de cambio de uso de suelo. \textbf{Desde enero de 2025, entonces, el desmonte de bosque
nativo está vedado en casi todo el departamento}, y a diferencia de los otros tres límites éste
no depende de una pendiente, de un ancho de calzada ni de una decisión futura: sale de una
cuenta hecha sobre la historia productiva de la cuenca.

\textbf{Puesto junto, eso describe un techo.} El valle está fraccionado, la ladera está ocupada
y además protegida, el bosque de la cuenca no es categorizable, y la ruta no puede crecer. \textbf{El dispositivo que produjo este
territorio está llegando al límite de lo que puede producir sin decisión.}

\subsection*{Y una ventana que está abierta ahora}

\begin{ficha}{Plan de Desarrollo Estratégico Salta 2050}
El \textbf{18 de mayo de 2026} se lanzó en el Centro de Convenciones de Salta el proceso
participativo de elaboración del \textbf{PDES 2050}, bajo el nombre «Salta 2050: Diálogos para
construir el futuro».

Financiamiento del \textbf{Consejo Federal de Inversiones}. Apertura a cargo del presidente del
Consejo Económico Social. Articulación a través de la Coordinación de Relaciones Políticas y
Planificación Estratégica. Jornada abierta al público \textbf{con cupo limitado} e inscripción
previa.

Se presenta como \textbf{actualización del PDES 2030}, creado en 2009 y actualizado en 2018.
\end{ficha}

\textbf{El plan que fijará el horizonte 2050 no existe todavía: se está escribiendo.}

En 2011 la Intendencia de La Caldera participó del Diagnóstico y el departamento apareció dos
veces en el documento final. \textbf{En 2026 el proceso empezó de nuevo, y todavía no está
decidido si esta vez aparece.}

\textbf{Y del lanzamiento se sabe algo más, que conviene poner al lado.} En un plenario de la
Cámara de Senadores provincial, representantes del Consejo Económico y Social expusieron la
puesta en marcha del plan y dejaron tres definiciones. La primera: que el plan busca adaptar las
metas a una provincia que ---según esa exposición--- transformó su matriz económica hacia
\textbf{la minería, la agroindustria y el turismo}. La segunda: que incorporará prospectiva por
escenarios y una \textbf{visión regionalizada}. Y la tercera: que el Consejo define metas,
mientras que el «cómo hacerlo» sigue siendo facultad exclusiva de los poderes públicos, y que
sus dictámenes técnicos son \textbf{no vinculantes}.

\begin{cautela}
Las tres definiciones alcanzan a este departamento, y en direcciones distintas.

\textbf{Las tres matrices que el plan declara son, justamente, las tres que este capítulo acaba
de descartar para el corredor.} La minería de La Caldera son veinte concesiones de áridos que
ocupan a cincuenta y una personas; el turismo, doscientas dieciocho; la agroindustria, ciento
noventa y tres sobre una superficie declarada que cayó. Si el plan 2050 se organiza sobre esas
tres matrices, el corredor vuelve a quedar sin casillero por la misma razón por la que no lo
tuvo en 2030, y no es exclusión: \textbf{lo que este departamento produce ---agua, piedra y
dormitorios--- no es una matriz productiva en el sentido en que esos planes usan la palabra.}

\textbf{La visión regionalizada es la única puerta.} Es el primer compromiso metodológico, en
los diecisiete años de planificación estratégica provincial que este libro repasa, que obliga a
bajar del agregado provincial al territorio. El reproche de este libro al PDES 2030 fue
exactamente ése: que sabe territorializar cuando quiere y que sobre este corredor no lo hizo.

\textbf{Y el carácter no vinculante marca la distancia con 2030.} Aquel plan se aprobó por
decreto invocando el artículo 77 de la Constitución provincial, que le daba «carácter
definitorio», y aun así este libro no encontró un solo acto del corredor que lo invoque
(capítulo \ref{cap:pdua}). Si lo que viene son dictámenes no vinculantes, la distancia entre el
plan y el expediente no se acorta: queda declarada de entrada.
\end{cautela}

\textbf{Y hay una continuidad que corresponde registrar, sin imputarle nada.} El PDES 2030 nació
en 2009 como iniciativa de la \textbf{Fundación Salta} (capítulo \ref{cap:pdua}). Diecisiete años
después, esa misma fundación integra el Consejo Económico y Social que impulsa la actualización,
y quien preside el Consejo es su vicepresidente. \textbf{No se sigue de eso ninguna imputación}:
se sigue que el órgano que fija el horizonte de largo plazo de la provincia tiene la misma
continuidad de personas e instituciones que este libro documenta en escalas mucho más chicas, y
que esa continuidad es, ella misma, un rasgo del objeto.

\begin{cautela}
\textbf{Toda esta subsección se apoya en partes de prensa y no en instrumentos.} Los partes del
Gobierno de la provincia del 12 y del 18 de mayo de 2026, el parte de la Cámara de Senadores
sobre el plenario ---que no consigna su fecha---, y las crónicas de los participantes
institucionales, consultados el 21 de septiembre de 2026. \textbf{No hay todavía acto
administrativo que cree el proceso}, ni términos de referencia publicados, ni cronograma. Lo que
se afirma aquí es que esas declaraciones se hicieron, no que el plan vaya a decir eso.
\end{cautela}

\begin{cautela}
\textbf{Ésa es la única afirmación de este capítulo que no es una proyección sino una
constatación:} hay un proceso de planificación abierto, con inscripción pública, y el
departamento tiene los materiales para presentarse a él. Este libro es uno de ellos; el apéndice
\ref{cap:propuestas}, otro.

\textbf{Lo que este libro no puede decir es si servirá de algo.} Un proceso participativo con
cupo limitado, en la capital, en día hábil, tiene las mismas barreras de acceso que la audiencia
pública del capítulo \ref{cap:vaqueros} y que el expediente en horario de oficina del capítulo
\ref{cap:ribera}. \textbf{El dispositivo también regula quién puede participar en las decisiones
que lo modificarían}, y eso es, exactamente, lo que la palabra nombra.
\end{cautela}

\section{Lo que refutaría todo este capítulo}

\textbf{Tres hallazgos lo desarmarían, y los tres son posibles.}

\textbf{Uno.} Que aparezca un acto de gobierno ---un decreto, un convenio, un acta--- que asigne al corredor, en el período en que se formó como periferia metropolitana, un perfil productivo que se haya sostenido. Si existe, hay estratega y el término no corresponde. \textbf{Y hay un tercer acto, incorporado al cierre, que conviene poner a prueba acá porque es el
que más se acerca}: el decreto 1.660 del 24 de abril de 1918, que resuelve que La Caldera
«\textit{no está en condiciones de ser administrado por una Municipalidad electiva}» y la
devuelve a Comisión Municipal (capítulo \ref{cap:siglo}). Es un acto de gobierno, fechado,
firmado y fundado, que decide sobre la condición institucional del departamento, y es la pieza
más parecida a un estratega que este relevamiento encontró. \textbf{No cumple la condición, y
conviene decir por qué}: no le asigna al corredor ningún perfil ni ningún destino ---se limita a
aplicar un umbral de renta a una planilla contable---, y cuatro semanas después el decreto 50 lo
extendió a los catorce departamentos por igual. Un criterio general aplicado con una cuenta no es
una decisión sobre este territorio. \textbf{Pero es el candidato más serio que el libro tiene, y
se deja consignado como tal para que se lo pueda discutir.} \textbf{El decreto 510 del 18 de abril
de 1922, que le devuelve al departamento su Comisión Municipal, es el candidato simétrico y
tampoco cumple}: decide sobre la existencia misma del órgano local ---que es más que decidir su
categoría---, pero lo hace en un acto general que alcanza a todos los departamentos «que en la
actualidad carecen de ella» y no le asigna al corredor ningún perfil. \textbf{Los dos actos
deciden qué gobierno tiene este territorio y ninguno de los dos dice para qué sirve}, que es
exactamente la diferencia entre un umbral y una estrategia. Los dos que este libro encontró ---el olivo en 1934 y el tabaco reconocido en 1980--- no cumplen esa condición: ninguno de esos perfiles se sostuvo, lo que es compatible con la ausencia de estratega. \textbf{Hay un tercer candidato, y es el que más cerca está de refutar este capítulo}: que el departamento haya sido pensado, y no sólo usado, como lugar de veraneo de la clase alta de la capital. El capítulo \ref{cap:poblacion} reúne lo que el archivo da a favor ---las casas quinta de 1927 y 1929, las truchas sembradas en 1940 y 1942, la iglesia ampliada en 1935, los apellidos de gobierno en el revalúo de 1944--- y lo que da en contra ---el olivo, el tabaco, el camino eliminado del plan en 1944, el puente postergado en 1935---, y concluye que el uso fue el más antiguo pero no fue asignado. \textbf{Si el padrón catastral de 1944 mostrara que sus titulares domiciliaban mayoritariamente en la capital, esa conclusión habría que revisarla, y con ella esta condición.}

\textbf{Dos.} Que el procesamiento censal de lugar de trabajo muestre que la mayoría de los
ocupados del departamento trabaja \textbf{dentro} de él. Si así fuera, la lectura residencial se
cae y con ella toda la sección prospectiva.

\textbf{Tres.} Que los organismos, requeridos, produzcan los documentos que este libro no
encontró: el aforo, el balance hídrico, el padrón de concesiones, el atlas del Código. Si todo
eso existe y simplemente no se publicó, \textbf{el problema es de publicidad y no de diseño
institucional}, y buena parte del apéndice \ref{cap:propuestas} sobra.

\textbf{Los tres se resuelven con los pedidos del apéndice \ref{cap:pedidos}}, ninguno de los
cuales había sido presentado al cerrarse este manuscrito.
